\documentclass[11pt]{article}
\usepackage{amsmath,amssymb,amsthm,mathtools,hyperref}

\begin{document}
\section*{The uncultivated microbial majority}
\label{prob:313}
\label{rank:201}
\par\smallskip

\noindent\textbf{Domain.} Uncultivated microorganisms; phenotype prediction. \quad \textbf{Retained status/source context.} Open genome-to-growth prediction programme. Sequence characterization and cultivation are distinct achievements; sequence detection does not establish autonomous growth conditions.

\paragraph{Definitions and assumptions.} Fix a declared set of nonpathogenic environmental taxa represented by quality-controlled genomes $g$. Let $e\in\mathcal E$ be a finite, normalized descriptor of environments already represented in observational data, and $c$ a community-context descriptor. Define $Y=1$ as sustained taxon-specific population increase over a prescribed horizon $H$, distinguished from DNA persistence or transient viability. A restricted model $p_\theta(g,e,c)=\mathbb{P}(Y=1\mid g,e,c)$ uses a fixed genomic feature map and a stated detection/error law. Unobserved environments are outside the validation domain unless explicitly modeled.

\paragraph{Problem statement.} Determine the identifiable growth-compatible region
    \[\mathcal V_g=\{(e,c):p_\theta(g,e,c)\geq p_0\},\]
    where $p_0\in(0,1)$ is prespecified. Quantify uncertainty across compatible parameter sets and require held-out taxon calibration within tolerance $\varepsilon$. A successful answer may be a bound or evidence that genome features alone are insufficient.

\paragraph{Additional binding assumptions and resolution criterion.} The region is a supported predictive growth region, not automatically the causal effect of setting environmental coordinates. Define community composition, absolute taxon abundance and the detection threshold jointly. Genome-feature transfer to withheld taxa requires a declared structural or phylogenetic model and a target taxon distribution. Region membership should be reported as inner/outer confidence sets where parameter ambiguity crosses the probability threshold. More explicitly, write $\mathcal V_{g,\theta}$ for the displayed region restricted to the declared joint prediction domain $D_g$ of environment--community pairs for genome $g$. For a nonempty compatible parameter set $C$, define
\[
\mathcal V_g^-(C)=\{(e,c)\in D_g:\text{for every }\theta\in C,\ p_\theta(g,e,c)\geq p_0\},\qquad
\mathcal V_g^+(C)=\{(e,c)\in D_g:\text{there exists }\theta\in C,\ p_\theta(g,e,c)\geq p_0\}.
\]
At population level, $C$ is the observational equivalence class under the declared observation designs and law; threshold membership is identified where all parameters in that class agree on membership. At finite sample size, $C=C(\mathsf D)$ is a data-dependent compatible set with a prespecified construction. To call the resulting sets confidence sets, declare a valid simultaneous coverage guarantee, its sampling law and its domain, including any claimed transfer to withheld taxa. An empty compatible set signals model rejection, not a certified growth region. Equality at $p_0$ is included; the existential outer-membership condition is not replaced by a possibly unattained supremum.

\paragraph{Declarations still required for a concrete instance.} The scheme does not select a taxonomic unit or genome-to-population mapping, genome quality criteria, initial-state law, abundance and time units, or a measurable sustained-growth event over $H$. These must be declared together with the observation law and treatment of detection limits; no particular cultivation or transfer protocol is implied. Declare the descriptor spaces and normalization, including whether ``finite'' means finite-valued or finite-dimensional, the joint prediction domain and reference law, and whether region claims are pointwise or modulo null sets. Pointwise claims require specified conditional-probability versions or model restrictions that determine them on the claimed domain. Observing an environmental descriptor alone does not establish support for every joint genome--environment--community combination. Declare when community and other predictors are available, which unmeasured states are averaged over, and the sampling and selection law. Prespecify the held-out taxon and context distribution, calibration functional, positive tolerance $\varepsilon$, holdout unit, separation of fitting and validation information, uncertainty guarantee and informative-output criterion. Evidence of genome-feature insufficiency must state its model, predictor-class or empirical scope; failure of one fitted model alone does not establish intrinsic insufficiency. These declarations retain the joint requirements of region inference, uncertainty quantification and held-out taxon calibration; the sources below do not select the outstanding conventions.

\paragraph{Alternative formulations.} (i) An autonomous-growth model fixes $c$ to an isolated-taxon context. (ii) A community-dependent model permits metabolic complementarity and targets abundance change in mixed populations. (iii) A phylogenetic transfer model predicts unfamiliar taxa under a declared distance limit; it is stronger than interpolation within already observed strains. These targets need not yield the same viable region.

\paragraph{Obstructions and partial progress.} Single-cell genomic studies greatly expanded the observed phylogenetic and coding diversity of uncultivated organisms. Such sequence progress does not determine all growth phenotypes. Missing genes, genome incompleteness, dormant states and community dependence produce ambiguity between intrinsic incapacity and unsuitable observation contexts. Selection into available growth datasets can bias apparent genomic predictors. The residual mathematical task is calibrated phenotype inference with a defined domain and observation law, not an unrestricted guarantee of conditions for every genome.

\paragraph{Sources (supplied provenance; not all rechecked).} \begin{itemize}
\item Supplied Wikipedia background: microbial dark matter. \url{https://en.wikipedia.org/wiki/Microbial_dark_matter} (Retained for provenance; not used as scientific verification of current open status.)
\item BMC Biology, Open questions in biology. \url{https://link.springer.com/collections/eiiddehjgj} (Collection index checked; historical research agenda, not direct verification of this particular statement.)
\item Rinke et al., Insights into the phylogeny and coding potential of microbial dark matter (2013). \url{https://pubmed.ncbi.nlm.nih.gov/23851394/} (Primary single-cell genomics checked; characterization rather than universal cultivation prediction.)
\end{itemize}


\section*{Binding catalogue conventions}
This edition revises both supplied catalogues and preserves original problem numbers 1--330. It adds one globally unique label \verb|\label{rank:x}| per problem, where smaller $x$ means greater provisional difficulty. A rank label is an anchor at that problem's section: \verb|\ref{rank:x}| returns the section number, not the literal difficulty number. The visible \texttt{rank:x} text and the index give the difficulty number itself.

The entries are \emph{parameterized mathematical research schemes}, including empirical prediction, causal identification, historical reconstruction, theory and engineering. They are not 330 certified unsolved mathematical conjectures. A rigorous instance fixes the objects and restrictions below; changing those objects changes the question. Established identities, finite calculations, non-identification constructions and solved special cases are not being posed anew as open. Some residual tasks concern informative bounds, biological realization or external validation rather than a new theorem.

Every main definition and problem statement received a mathematical editorial review. Each entry now has an additional binding-conditions paragraph; direct corrections are separately recorded in the accompanying audit. The alternative formulations inherit the applicable conventions, but are different targets unless equivalence is proved. Where an old broad phrase conflicts with an explicit restriction below it, the explicit restriction controls.

\textbf{Literature boundary.} Thirteen selective primary-literature checks cover eighteen entries in this revision; their access levels and scope are recorded explicitly. The supplied bibliography and its earlier source-access descriptions are retained as provenance, not represented as newly verified. Neither this review nor a source's description proves the absence of later solutions. The date of this edition is 13 October 2026; there is no claim of exhaustive current-status certification for all 330 entries.

\subsection*{Binding conventions for a concrete instance}
\paragraph{Objects, quantifiers and non-vacuity.}
An instance declares measurable state, input, output and parameter spaces; a fixed candidate family; admissible initial laws, interventions and environments; a finite horizon or specified asymptotic regime; observation and sampling laws; and performance metrics, positive normalization scales and positive tolerances. These are chosen before fitting or testing, not adjusted after seeing failures. Shared parameters must genuinely be shared across the declared validation domain. Finite dimension alone does not restrict arbitrary real-number encodings: specify a finite library, bounded description length, or regularity/complexity bounds. Exclude constant or no-signal solutions through a fixed informative-output or baseline-improvement requirement where appropriate. A universally quantified result cannot be established by choosing one trivial favorable instance.

\paragraph{Probability laws and data.}
Every expectation or probability refers to the specified stochastic law, including initial uncertainty, process noise, observation noise and sampling as appropriate. A deterministic equation is not a probability model without a declared random input or initial law. ``Observed'' targets in population-level formulas mean population observables $g_j^*$, not noisy estimates asserted to be exact. Let $\widetilde g_j$ be independent validation estimates and suppose
\[
\mathbb{P}\{\lvert\widetilde g_j-g_j^*\rvert\le r_j\text{ for every }j\}\ge1-\alpha.
\]
Then an uncertainty-aware sufficient certification rule for positive target tolerances $\varepsilon_j$ is
\[
\max_j\frac{|\widehat g_j-\widetilde g_j|+r_j}{\varepsilon_j}\le1.
\]
On that simultaneous confidence event it implies every population error is at most its tolerance, by the triangle inequality. Vector/norm versions use corresponding valid simultaneous bounds. Training uncertainty must also be included when the claim averages over repeated training data. Evaluation on finitely many conditions does not establish a continuous-domain supremum without covering and regularity bounds.

\paragraph{Dynamics and well-posedness.}
For ODEs and SDEs give coefficient domains, initial laws and existence/uniqueness conditions in the intended solution sense, or keep the entire admitted solution set. For stochastic jump models specify nonnegative off-diagonal rates, diagonal normalization, rate orientation and non-explosion. For PDEs specify spatial domains, boundary/initial conditions, conservation or source terms, and a weak or classical solution class. Moving-domain concentration equations include dilution/transport terms. A stationary law is used only for a time-homogeneous augmented model under existence assumptions; a stationary approximation to a driven system must be justified. Absorbing extinction states are retained. Survival-conditioned, quasi-stationary and unconditional laws are distinct and must not be substituted for one another.

\paragraph{Identification and causal meaning.}
Two parameter values are observationally equivalent when they induce the same full distribution under all allowed observation/intervention designs, not merely the same likelihood value for one sample. A target is identified only if constant on the equivalence class. Gauge symmetries, label switching and reference conventions are quotiented or fixed. Finite-data compatible sets use a declared confidence, likelihood or discrepancy rule; an empty set signals model rejection. Bayesian credible sets are not automatically confidence sets. A prior or regularizer does not create data identification. Bounds intended to resolve a question must be informative at a predeclared width or decision threshold.

A causal contrast uses well-defined intervention versions and a stated causal model. Random assignment, or explicit consistency, overlap and exchangeability assumptions, is required for causal interpretation from observations. Interference and time-varying confounding are represented when present. Fixing a mediator by intervention differs from conditioning on its observed value. Marginal potential-outcome laws do not determine individual cross-world effects or an arbitrary joint coupling. Ethically admissible designs, rather than unconstrained interventions, define the experimental menu.

\paragraph{Metrics, units, optimization and limits.}
Specify the sampling or reference measure for densities and integrals, the log base, and dimensionless arguments of logarithms. Total variation is $\sup_A|P(A)-Q(A)|$; Wasserstein distances require the declared ground metric and appropriate moment conditions; KL divergence uses its stated absolute-continuity convention. All normed readouts are scaled consistently, and error probabilities use $0<\alpha<1$. A ratio has a positive denominator or an explicit undefined/zero convention. Thermodynamic entropy production is nonnegative in expectation under the relevant hypotheses; individual stochastic realizations need not be. An infimum is not an attained minimum without compactness/coercivity and continuity or another existence argument. Fix the order of large-population, long-time and vanishing-mutation limits. Stability claims declare perturbation space and symmetries.

\subsection*{How the provisional difficulty order was assigned}
Four editorial components are scored from 1 (lower burden) to 5 (higher burden): (a) formal or inferential difficulty, (b) access to evidence and identification burden, (c) mechanistic/computational integration, and (d) independent validation and transfer. Their unweighted total orders the entries. For selected entries in the highest bands, an explicit editorial priority breaks equal totals; all other ties are broken by original problem number. Equal totals are \emph{not} evidence of reliably different true difficulties. Scores are judgments, not measurements, completion-time estimates or consensus rankings.

The object ranked is an informative, nondegenerate resolution of the \emph{scoped} scientific target, with the stated budget and validation domain. A sharp useful impossibility result can qualify; an already known identity, vacuous bound, arbitrary favorable model or lookup table cannot. Rank 330 is merely the least difficult relative assessment in this catalogue, not a claim of an easy biological solution. Broad historical titles do not automatically outrank precisely formulated mechanistic tasks.



\end{document}
